\documentclass[11pt]{article}
\usepackage{fontspec}
\usepackage[french,english]{babel}
\usepackage{amsmath}
\usepackage{amssymb}
\usepackage{amsthm}
\usepackage{mathrsfs}
\usepackage{tikz-cd}
\usepackage{enumitem}
\usetikzlibrary{babel}

\newcommand{\printedpage}[1]{\marginpar{\footnotesize p.~#1}}
\newcommand{\msp}[1]{\marginpar{\footnotesize [#1]}}

% footnote marks as printed: "(1)" in the text, "1" before the footnote text
\makeatletter
\renewcommand\@makefnmark{\hbox{\textsuperscript{\normalfont(\@thefnmark)}}}
\renewcommand\@makefntext[1]{\parindent 1em\noindent\hb@xt@1.8em{\hss\textsuperscript{\normalfont\@thefnmark}}#1}
\makeatother

\begin{document}

% ===================== PDF page 1 =====================
\printedpage{}% no page number printed on this page
% head: (none)

\begin{center}
\vspace*{6cm}

{\LARGE \textsf{THE DERIVATORS}}

\vspace{1cm}

{\large \textbf{ALEXANDRE GROTHENDIECK}}

\vspace{2cm}

{\large \textbf{Chapter XIV}}

\vspace{0.3cm}

{\large \textbf{h-cartesian and h-cocartesian squares}}
\end{center}

\clearpage

% ===================== PDF page 2 =====================
\printedpage{}% no page number printed on this page
% head: (none)

\vspace*{6cm}

\noindent This text has been deciphered and transcribed in \LaTeXe{} by M. Künzer. It has been edited by M. Künzer, J. Malgoire, and G. Maltsiniotis. The transcription is as faithful as possible to the manuscript. For the few obvious corrections, or rare comments by the editors, as well as for the original numbering of the manuscript pages, typewriter characters \texttt{[typewriter]} in square brackets are used. A question mark in square brackets means that one is not sure of the preceding word.

\noindent This chapter is unfinished. In the first section, Grothendieck introduces a notion of left (resp. right) ``quadrangulated category''. The text breaks off abruptly at the beginning of the second section, devoted to examples, when Grothendieck realizes, probably, that the main example he has in mind (the homotopy category of a category of models equipped with homotopically cartesian (resp. cocartesian) squares) does \textit{not} satisfy the axiomatics of quadrangulated categories. This is briefly explained in an afterword by the editors.

\noindent This edition is provisional. Remarks, comments and corrections are welcome. Send a message to:

\noindent \textbf{[e-mail address removed]}

\hfill G. Maltsiniotis\hspace*{2cm}

\clearpage

% ===================== PDF page 3 =====================
\printedpage{1}
% head: A. Grothendieck : Dérivateurs, Ch. XIV, Édition des Universités de Montpellier II et Paris VII
% foot: Version 7 octobre 2010

\msp{page 1}

\section*{1\quad Axioms for quadrangulated structures (an aborted attempt)}

Let $(\mathcal{M},W)$ be a category of models admitting a theory of the homotopy fiber. Then we have in $\mathrm{Ho}_W=\mathcal{M}W^{-1}$ a notion of \textit{$W$-cartesian square} (or $W$-homotopically cartesian), by which we mean a square diagram
\[
\begin{tikzcd}
X_1 \arrow[d] & X_3 \arrow[l] \arrow[d] \\
X_0 & X_2 \arrow[l]
\end{tikzcd}
\]
in $\mathrm{Ho}_W$ coming, up to isomorphism, from a $W$-cartesian square in $\mathcal{M}$ – thus isomorphic, in $\underline{\mathrm{Hom}}(\boldsymbol{\Delta}_1\times\boldsymbol{\Delta}_1,\mathrm{Ho}_W)$, to a \textit{cartesian} square \texttt{[in $\mathcal{M}$]} % [sic] a un carré -> à un carré, preposition printed without its grave accent
\[
\begin{tikzcd}
Y_1 \arrow[d, "p_1"'] & Y_3\ (\simeq Y_1\times_{Y_0}Y_2) \arrow[l] \arrow[d] & \\
Y_0 & Y_2 \arrow[l, "p_2"] & ,
\end{tikzcd}
\]
with $p_1,\ p_2\in\mathrm{Fib}_W$. (In fact, it suffices to have $p_1$ or $p_2$ in $\mathrm{Fib}_W$ to ensure the $W$-cartesian character of the cartesian square under consideration.) I would like to axiomatize this situation, mainly to make a list of the essential useful properties of this notion, and of the dual notion of $W$-cocartesian square (or $W$-homotopically cocartesian), in the case where $(\mathcal{M},W)$ admits a theory of the homotopy cofiber. The species of structure, namely that of ``quadrangulated categories'' (left, right, two-sided) is a variant, in a \msp{page 2} nonadditive case, of the notion of triangulated category isolated by \textsc{Verdier}. Like the latter, this notion is in fact too weak – it is \textit{not} ``the right one'' – but it is of provisional practical interest. The ``right'' notion (if there is one apart from that of derivator) should work with a finer notion of square and of $W$-cartesian square, working systematically in $\mathrm{HO}_W(\boldsymbol{\Delta}_1\times\boldsymbol{\Delta}_1)=\mathcal{M}(\boldsymbol{\Delta}_1\times\boldsymbol{\Delta}_1)W(\boldsymbol{\Delta}_1\times\boldsymbol{\Delta}_1)^{-1}$ (where the ``real'' $W$-cartesian squares live!), and similar categories, such as $\mathrm{HO}_W(\Psi)$, $\mathrm{HO}_W(\Phi)$. But to include them in an axiomatics making no mention of a category of models $(\mathcal{M},W)$, manifestly amounts to putting one's finger in the gears of the notion of derivator, and that is not my purpose for the moment.

\noindent\textbf{Definition 1.} I call a \textit{``category with h-cartesian squares'',} or a \textit{``left quadrangulated category'',} a category $\mathcal{H}$, equipped with a subset Cart of the set $\mathrm{Hom}(\boldsymbol{\Delta}_1\times\boldsymbol{\Delta}_1,\mathcal{H})$
\printedpage{2}
% head: A. Grothendieck : Dérivateurs, Ch. XIV, Édition des Universités de Montpellier II et Paris VII
% foot: Version 7 octobre 2010
of square diagrams \texttt{[by square we mean commutative square]}
\[
\begin{tikzcd}
X_1 \arrow[d] & X_3 \arrow[l] \arrow[d] \\
X_0 & X_2 \arrow[l]
\end{tikzcd}
\]
in $\mathcal{H}$, squares called \textit{``h-cartesian''} (h as in ``homotopically''), this \msp{page 3} datum being subject to the following axioms, Cart (1) to Cart (\ldots).

\noindent Since I am not sure I know all the essential properties in advance, the number of these axioms remains open for the moment.

\noindent\textbf{Cart (1) (Symmetry).} If a square
\[
\begin{tikzcd}
X_1 \arrow[d] & X_3 \arrow[l] \arrow[d] \\
X_0 & X_2 \arrow[l]
\end{tikzcd}
\]
is h-cartesian, likewise the symmetric square (obtained by interchanging the roles of $X_1$ and $X_2$).

\noindent\textbf{Cart (2) (Trivial squares).} Let
\[
\begin{tikzcd}
X_1 \arrow[d, "p_1"'] & X_3 \arrow[l, "q_2"'] \arrow[d, "q_1"] \\
X_0 & X_2 \arrow[l, "p_2"]
\end{tikzcd}
\]
be a square in $\mathcal{H}$, with $p_1$ an isomorphism. Then the square is h-cartesian if and only if $q_1$ is an isomorphism.

\noindent\textbf{Cart (3) (Composite squares).} Let there be a commutative diagram
\[
\begin{tikzcd}
X'_0 \arrow[d] \arrow[dr, phantom, "\mathrm{I}"] & X'_1 \arrow[l] \arrow[d] \arrow[dr, phantom, "\mathrm{II}"] & X'_2 \arrow[l] \arrow[d] & \\
X_0 & X_1 \arrow[l] & X_2 \arrow[l] & ,
\end{tikzcd}
\]
whence the squares I, II and III:
\[
\begin{tikzcd}
X'_0 \arrow[d] & X'_2 \arrow[l] \arrow[d] \\
X_0 & X_2 \arrow[l]
\end{tikzcd}
\]
\printedpage{3}
% head: A. Grothendieck : Dérivateurs, Ch. XIV, Édition des Universités de Montpellier II et Paris VII
% foot: Version 7 octobre 2010
(which we call the \textit{composite square} (``horizontal'') of the squares I and II). Suppose I h-cartesian. Then II is h-cartesian if and only if III is \texttt{[see editors' afterword]}.

\noindent\textbf{Corollary.} \textit{Let there be a commutative cubical diagram, which we can regard as a homomorphism from one square} (\textit{upper face of the cube}) \textit{to another} (\textit{lower face}):
\[
\begin{tikzcd}
 & {} \arrow[dl] \arrow[dd] & & {} \arrow[ll] \arrow[dl] \arrow[dd] \\
{} \arrow[dd] & & {} \arrow[ll, crossing over] & \\
 & {} \arrow[dl] & & {} \arrow[ll] \arrow[dl] \\
{} & & {} \arrow[ll] \arrow[from=uu, crossing over] & .
\end{tikzcd}
\]
\textit{If all the faces of the cube}\footnote{In fact, the lower face, and the two parallel lateral faces (e.g. the front one and the back one) suffice.} \textit{except the upper face are} \msp{page 4} \textit{h-cartesian squares (namely the lower face and the four lateral faces), the same holds for the upper face.}\footnote{\textbf{Corollary} (of Cart (3) and Cart (2)): \textit{every square isomorphic to an h-cartesian square is h-cartesian.}}

\noindent These axioms are in the nature of stability axioms, giving us criteria asserting that certain squares are h-cartesian, if such and such conditions are fulfilled (e.g. that certain other squares are already h-cartesian). I have the impression of having exhausted here the relevant properties of this type (\textit{i.e.} that the others should be deducible from them). At present, I am dealing with properties of existence and uniqueness of h-cartesian squares. % [sic] carrés h-cartésien -> carrés h-cartésiens, missing plural agreement

\noindent\textbf{Cart (4) (Fiber product axiom).} Every diagram
\[
\begin{tikzcd}
X_1 \arrow[d] & \\
X_0 & X_2 \arrow[l]
\end{tikzcd}
\]
in $\mathcal{H}$ can be completed to an h-cartesian square
\[
\begin{tikzcd}
X_1 \arrow[d] & X_3 \arrow[l] \arrow[d] & \\
X_0 & X_2 \arrow[l] & .
\end{tikzcd}
\]
The following axiom (\textit{cf.} its corollary) will ensure among other things that the h-cartesian square in question is \textit{unique up to isomorphism} – but one should note that it is in general not unique ``up to unique isomorphism''.

\end{document}
